\SourceLeaf{308}
\fancyhead[L]{\small UNTERSUCHUNGEN ÜBER GEGENSTÄNDE DER HÖHERN GEODAESIE}
\renewcommand{\Article}[1]{\par\medskip\Needspace{4\baselineskip}\pdfbookmark[1]{Art. #1}{geodaesie2-art-#1}\centerline{\textbf{#1.}}\nobreak\smallskip}
\begin{center}
UNTERSUCHUNGEN\par\bigskip
ÜBER\par\bigskip
{\large GEGENSTÄNDE DER HÖHERN GEODAESIE.}\par\medskip
\rule{.65in}{.3pt}
\end{center}
\noindent\hspace*{\parindent}Die Aufgabe, aus der Grösse der Seite eines Dreiecks auf der Erdoberfläche, dem Azimuthe an dem einen Endpunkte, und der geographischen Breite dieses Endpunkts abzuleiten das Azimuth an dem andern Endpunkte, dessen Breite und den Längenunterschied beider Punkte, gehört zu den Hauptgeschäften der höhern Geodäsie. Für den Fall der Kugelfläche ist der Zusammenhang zwischen jenen sechs Grössen am Schluss der ersten Abhandlung in der einfachsten und zur schärfsten Rechnung geeigneten Form aufgestellt, welche auch leicht zu einer bequemen Auflösung der Aufgabe selbst benutzt werden kann. Es wird dadurch das Verlangen nach dem Besitz einer analogen unmittelbar für die Ellipsoidfläche gültigen Auflösungsart erweckt, und der Zweck der gegenwärtigen Abhandlung ist, eine solche zu entwickeln. Vorher soll jedoch erst die Auflösung für den Fall der Kugelfläche in ein noch helleres Licht gestellt werden. Des bequemern Zurückweisens wegen lasse ich die Zahlenbezeichnung der Artikel sich an die erste Abhandlung anschliessen.

\Article{18}
Um den Grad der Genauigkeit, welcher durch die Formeln des 17. Art. erreicht wird, besser beurtheilen zu können, werden noch die Glieder der nächstfolgenden Ordnung entwickelt werden müssen; es ist jedoch wohl der Mühe
